\section{Complete finite-value findings and the remaining contribution question}
\label{sec:finite-value-findings}

I retain all 819 emitted architecture profiles under the same 273 certified
financial prefixes. Range, clipped and dyadic estimators share the actual
finite-output support and moments; the unit-range flag also receives the
strengthened whole-mean prior. The table selects over the same finite pool and
two proved scalar controllers at per-price failure probability $0.01$.
Every selected implementation retains the original deterministic allowance
at most $0.005$ and estimation allowance $0.005/H$.

\begin{center}
\begin{tabular}{lrrrr}
\toprule
Family & Lower T cap & Equal & Higher & Minimum ratio\\
\midrule
Range flag &16&11&0&0.352221\\
Clipped flag &17&10&0&0.179376\\
Best single flag &17&10&0&0.179376\\
Dyadic values &0&0&27&10.312280\\
\bottomrule
\end{tabular}
\end{center}

Ratios in this table divide by the selected unit-range flag's complete
logical T cap. The strongest single-flag construction reduces 17 of 27
contract budgets and leaves ten unchanged, with maximum reduction approximately
$82.1$ per cent. Against that stronger single flag, the selected dyadic
construction loses every contract: its T cap is between approximately
$21.00$ and $122.03$ times the opponent's cap. This is an adverse finding for
the implemented global-moment dyadic family. It is not a theorem excluding
all variance-aware quantum means, locally informed buckets, better allocations
or other coherent encodings. Nor does a comparison of upper caps establish
realised stochastic-cost dominance or physical quantum advantage.

The full archive \texttt{results/variance\_oracle\_comparison.json.gz}
retains every accepted and rejected scalar candidate, 42,024 bucket records,
resource charges and 162 contract/family minima. An independent auditor checks
all 819 rows, 8,892 scalar candidates, 3,276 dyadic candidates, 8,190 profile
minima, the final selections and the summary arithmetic. The source and
consultation records remain separate from these mathematical and numerical
findings. The earlier 58 protected numerical outputs are unchanged.

The positive project-specific consequence is the complete budget reduction
that a checked support/moment/numerical-transfer certificate supplies to a
competent single flag. The finite circuit and resource comparison make that
consequence constructive. The generic ingredients have established
predecessors, and I have not established that this exact interface argument is
absent from prior literature. Clipping moment bounds, dyadic mean estimation,
range normalisation, geometric controls and conditional smoothing are not
claimed as inventions. The bounded-support comparison with \citet{liuli2009} remains
incomplete at its equation and proof level in the source ledger. Stronger
higher-dimensional classical pricing and compatible modern quantum encodings
remain necessary before a wider contribution or practical advantage claim.

The deterministic two-asset findings already defeat an inference based only
on beating certified Monte Carlo. I therefore retain the question at the
transferable financial certificate, encoding and complete cost interface,
rather than treating a different asset label or a larger list of models as a
contribution. A new conditional or bucket-local certificate would merit
further development only if it changes a full budget against these opponents
and supplies an original substantive argument supported by the literature.
